\section{Chern Numbers \& Pontrjagin Numbers}
\label{sec:ChernNumbers&PontrjaginNumbers}

\subsection{Chern Numbers}
\begin{definition}
    Let $K^n$ be a compact complex manifold, 
    $\dim_\mathbb{C}K=n$. 
    For each partition $I=i_1,\dots,i_r$ of $n$, 
    we define the $I$-th Chern number to be the integer 
    \begin{equation*}
        c_I[K^n]=c_{i_1}\cdots c_{i_r}[K^n]
        :=\langle c_{i_1}(\tau)\cdots c_{i_r}(\tau),\mu_{K^n}\rangle,
    \end{equation*}
    where $\mu_{K^n}$ is the fundamental homology class w.r.t. 
    the orientation determined by the complex structure.

    If $I$ is a partition of other number, we define $c_I[K^n]$ to be $0$. 
\end{definition}

\begin{example}
    For $K^n=\mathbb{C}P^n$, we have
    \begin{align*}
        c_I[\mathbb{C}P^n] 
        &= \langle c_{i_1}(\tau)\cdots c_{i_r}(\tau),
        \mu_{\mathbb{C}P^n}\rangle \\
        &=\left\langle\binom{n+1}{i_1}\cdots\binom{n+1}{i_r}a^n,
        \mu_{\mathbb{C}P^n}\right\rangle \\
        &=\binom{n+1}{i_1}\cdots\binom{n+1}{i_r}.
    \end{align*}
\end{example}

\subsection{Pontrjagin Numbers}
\begin{definition}
    Let $M^{4n}$ be a compact oriented smooth manifold of dimension $4n$. 
    For each partition $I=i_1,\dots,i_r$ of $n$, 
    we define the $I$-th Pontrjagin number to be the integer 
    \begin{equation*}
        p_I[M^{4n}]=p_{i_1}\cdots p_{i_r}[M^{4n}]
        :=\langle p_{i_1}(\tau)\cdots p_{i_r}(\tau),\mu_{M^{4n}}\rangle,
    \end{equation*}
    where $\mu_{M^{4n}}$ is the fundamental homology class w.r.t. 
    the given orientation.

    If $I$ is a partition of other number, 
    we define $p_I[M^{4n}]$ to be $0$.
\end{definition}

\begin{example}
    For $M^{4n}=\mathbb{C}P^{2n}$, we have
    \begin{align*}
        p_I[\mathbb{C}P^{2n}] 
        &= \langle p_{i_1}(\tau)\cdots p_{i_r}(\tau),
        \mu_{\mathbb{C}P^{2n}}\rangle \\
        &=\left\langle\binom{2n+1}{i_1}\cdots\binom{2n+1}{i_r}a^{2n},
        \mu_{\mathbb{C}P^{2n}}\right\rangle \\
        &=\binom{2n+1}{i_1}\cdots\binom{2n+1}{i_r}.
    \end{align*}
\end{example}

\noindent\textbf{Oriented reversing map.} 
If we reverse the orientation of $M^{4n}$, since the definition of
Pontrjagin classes do not depend on the orientation, 
we have the following
\begin{gather*}
    p_k(M)\mapsto p_k(M), \\
    \mu_M\mapsto -\mu_M, \\
    p_I[M]\mapsto -p_I[M].
\end{gather*}
Thus the sign of Pontrjagin numbers change when 
reversing the orientation of $M$. 
If $M$ has a nonzero Pontrjagin number, 
$M$ cannot possess any orientation reversing diffeomorphism. 

\noindent\textbf{Manifold as a boundary.} 
If $M^{4n}$ is a boundary of some smooth, 
compact, oriented $(4n+1)$-dimensional manifold-with-boundary. 
Then all Pontrjagin numbers of $M$ must be zero. 
Thus if some $p_I[M]\neq0$, 
then $M$ cannot be the boundary of such manifold.

\subsection{Symmetric Functions}
We know that all symmetric polynomials in 
$\mathbb{Z}[t_1,\dots,t_n]$ is itself a polynomial ring: 
\begin{equation*}
    \mathcal{S} = \mathbb{Z}[\sigma_1,\dots,\sigma_n]\subset
    \mathbb{Z}[t_1,\dots,t_n],
\end{equation*}
where $\sigma_i$ is the $i$-th elementary symmetric polynomial.
\begin{equation*}
    1+\sigma_1+\cdots+\sigma_n = (1+t_1)\cdots(1+t_n).
\end{equation*}
If we focus on the symmetric polynomials of degree $k$, 
denoted by $\mathcal{S}^k$, clearly 
it is a free abelian group with basis
\begin{equation*}
    \sigma_{i_1}\cdots\sigma_{i_r},\quad
    I=i_1,\dots,i_r\text{ is a partition of } k.
\end{equation*}
For some reason, it is more convenient to use another basis of 
$\mathcal{S}^k$, which is described as follows.

Define two monomials in $\mathbb{Z}[t_1,\dots,t_n]$ to be 
equivalent if they differ by a permutation of the variables.
We define 
\begin{equation*}
    \sum t_1^{a_1}\cdots t_r^{a_r}
\end{equation*}
to be the summation of all distinct monomials equivalent to
$t_1^{a_1}\cdots t_n^{a_n}$. Clearly, it is a symmetric polynomial.

\begin{lemma}
    \begin{equation*}
        \sum t_1^{a_1}\cdots t_r^{a_r}
    \end{equation*}
    where $a_1\cdots a_r$ ranges over all partitions 
    of $k$ with length $r\leq n$,
    form another basis of $\mathcal{S}^k$.
\end{lemma}

Now for each partition $I=i_1,\dots,i_r$ of $k$, 
we define a polynomial $s_I$ in $k$ variables as follows:
choose $n\geq k$, define 
\begin{equation*}
    s_I(\sigma_1,\dots,\sigma_k)
    :=\sum t_1^{i_1}\cdots t_r^{i_r},
\end{equation*}
both sides are symmetric polynomials in $n$ variables.

Notes that $s_I$ does not depend on the choice of $n$, 
for if we choose another $n'>n$, 
then we can set $t_{n+1}=\cdots=t_{n'}=0$ to get back to the previous case.

Clearly, all $s_I$ form a basis of $\mathcal{S}^k$ as $I$ ranges over
all partitions of $k$.

Now let's look at the cohomology ring 
$H^*(\mathrm{Gr}_n(\mathbb{C}^\infty);\mathbb{Z})$.
the splitting map 
\begin{equation*}
    \phi:\mathbb{C}P^\infty\times\cdots\times\mathbb{C}P^\infty
    \ra \mathrm{Gr}_n(\mathbb{C}^\infty)
\end{equation*}
induces an injective homomorphism of cohomology rings
\begin{equation*}
    \phi^*:H^*(\mathrm{Gr}_n(\mathbb{C}^\infty);\mathbb{Z})
    \ra H^*(\mathbb{C}P^\infty\times\cdots\times\mathbb{C}P^\infty;
    \mathbb{Z})
    \cong \mathbb{Z}[t_1,\dots,t_n],
\end{equation*}
where $t_i$ is the generator of the $i$-th factor.
$\phi^*$ maps 
$H^*(\mathrm{Gr}_n(\mathbb{C}^\infty);\mathbb{Z})=\mathbb{Z}[c_1,\dots,c_n]$ 
isomorphically onto the symmetric polynomial subring 
$\mathcal{S}\subset\mathbb{Z}[t_1,\dots,t_n]$, 
with $c_i$ mapped to $\sigma_i$.

Since $\phi^*s_I(c_1,\dots,c_k) = s_I(\sigma_1,\dots,\sigma_k)$ 
forms a basis of $\mathcal{S}^k$ as $I$ ranges over all partitions of $k$, 
by the injectivity of $\phi^*$, $s_I(c_1,\dots,c_k)$ also forms a basis of 
$H^{2k}(\mathrm{Gr}_n(\mathbb{C}^\infty);\mathbb{Z})$.

\subsection{A Product Formula}
Let $\omega$ be an $n$-dimensional complex vector bundle over
a compact complex manifold $B$. For each $k\geq0$ and partition 
$I=i_1,\dots,i_r$ of $k$, we define
\begin{equation*}
    s_I(c(\omega)) := s_I(c_1(\omega),\dots,c_k(\omega))\in
    H^{2k}(B;\mathbb{Z}).
\end{equation*}
sometimes we simply write it as $s_I(c)$. 

\begin{lemma}[Thom]
    \begin{equation*}
        s_I(c(\omega\oplus\omega')) =
        \sum_{J,K} s_J(c(\omega)) s_K(c(\omega')),
    \end{equation*}
    where the summation is taken over all partitions $J,K$
    with juxtapositon $JK$ equal to $I$.
\end{lemma}
\begin{remark}
    We can view this formula as a generalization of the Whitney sum formula.
\end{remark}
\begin{proof}
    Let $\phi:F(B)\to B$ be the flag manifold of $\omega\oplus\omega'$,
    with the splitting bundle
    \begin{equation*}
        \phi^*(\omega) = \eta_1\oplus\cdots\oplus\eta_n, \quad
        \phi^*(\omega') = \eta_{1}'\oplus\cdots\oplus\eta_{n'}'.
    \end{equation*}
    Then we have
    \begin{gather*}
        \phi^*c(\omega) = 
        (1+c_1(\eta_1))\cdots(1+c_1(\eta_n)) =: 
        (1+t_1)\cdots(1+t_n), \\
        \phi^*c(\omega') = 
        (1+c_1(\eta_1'))\cdots(1+c_1(\eta_{n'}')) =: 
        (1+t_{n+1})\cdots(1+t_{n+n'}).
    \end{gather*}
    Thus
    \begin{align*}
        \phi^*s_I(c(\omega\oplus\omega')) 
        &= \sum t_1^{i_1}\cdots t_r^{i_r} \\
        &= \sum_{
        \begin{subarray}{c}
            JK = I, \\ J = j_1,\dots,j_{r_1}, \\ K = k_1,\dots,k_{r_2}
        \end{subarray}
        } 
        \left(\sum t_1^{j_1}\cdots t_{r_1}^{j_{r_1}}\right)
        \left(\sum t_{n+1}^{k_1}\cdots t_{n+r_2}^{k_{r_2}}\right) \\
        &= \sum_{JK=I} s_J(c(\eta_1\oplus\cdots\oplus\eta_n)) s_K(c(\eta_1'\oplus\cdots\oplus\eta_n')) \\
        &= \sum_{JK=I} s_J(\phi^*c(\omega)) s_K(\phi^*c(\omega')) \\
        &= \phi^*\left(\sum_{JK=I} s_J(c(\omega)) s_K(c(\omega'))\right).
    \end{align*}
    Since $\phi^*$ is injective, we get the desired formula.
\end{proof}

Now for a compact complex manifold $K^n$. For each partition
$I=i_1,\dots,i_r$ of $n$, we denote
\begin{equation*}
    s_I[K^n] = s_I(c)[K^n] := \langle s_I(c(\tau)), \mu_{K^n}\rangle\in\mathbb{Z}.
\end{equation*}
This characteristic number is clearly a linear combination of
Chern numbers.

\begin{corollary}
    \begin{equation*}
        s_I[K^m\times L^n] =
        \sum_{JK=I} s_J[K^m] s_K[L^n].
    \end{equation*}
    where the summation is taken over all partitions $J$ of $m$ and $K$ of $n$
    with juxtapositon $JK$ equal to $I$.
\end{corollary}
\begin{proof}
    Since $\tau_{K\times L} = \tau_K \times \tau_L = \pi_1^*\tau_K \oplus \pi_2^*\tau_L$,
    where $\pi_1,\pi_2$ are the projections from $K\times L$ to
    $K$ and $L$ respectively, 
    \begin{align*}
        s_I[K\times L]
        &= \langle s_I(c(\tau_{K\times L})), \mu_{K\times L}\rangle \\
        &= \sum_{JK=I} \left\langle s_J(c(\pi_1^*\tau_K)) s_K(c(\pi_2^*\tau_L)),
        \mu_{K}\times \mu_{L} \right\rangle \\
        &= \sum_{JK=I} \langle s_J(c(\tau_K)), \mu_K \rangle
        \langle s_K(c(\tau_L)), \mu_L \rangle \\
        &= \sum_{JK=I} s_J[K] s_K[L].
    \end{align*}
\end{proof}

\begin{corollary}
    As a special case, $s_{m+n}[K^m\times L^n] = 0$, for $m,n>0$.
\end{corollary}

There are analogous results for Pontrjagin classes and numbers. 
Let $\xi$ be a real vector bundle over a compact smooth manifold $B$. 
For each $k\geq0$ and partition $I=i_1,\dots,i_r$ of $k$, we define
\begin{equation*}
    s_I(p(\xi)) := s_I(p_1(\xi),\dots,p_k(\xi))\in
    H^{4k}(B;\mathbb{Z}).
\end{equation*}
There are similar product formulas for Pontrjagin classes 
\begin{equation*}
    s_I(p(\xi\oplus\xi')) \equiv
    \sum_{J,K} s_J(p(\xi)) s_K(p(\xi')),
\end{equation*}
modulo elements of $2$-torsions. Since pairing with the fundamental class
annihilates all $2$-torsions, we have 
\begin{equation*}
    s_I[M\times N] = \sum_{JK=I} s_J[M] s_K[N],
\end{equation*}
Notes that these characteristic numbers
will be zero unless both the dimensions of $M$ and $N$ are divisible by $4$.

\subsection{Linear Independence of Chern Numbers and of Pontrjagin Numbers}
\begin{theorem}[Thom]
    Let $K^{1},\dots,K^{n}$ be compact complex manifolds
    with $s_k(c)[K^{k}]\neq0$. Then the $p(n)\times p(n)$ matrix
    \begin{equation*}
        \Big(c_{i_1}\cdots c_{i_r}[K^{j_1}\times\cdots\times K^{j_s}]\Big)
    \end{equation*}
    is nonsingular, where both $I=i_1,\dots,i_r$ and
    $J=j_1,\dots,j_s$ range over all partitions of $n$. 
\end{theorem} 
Similarly, we have the following theorem for Pontrjagin numbers.
\begin{theorem}[Thom]
    Let $M^{4},\dots,M^{4n}$ be compact oriented smooth manifolds
    with $s_k(p)[M^{4k}]\neq0$. Then the $p(n)\times p(n)$ matrix
    \begin{equation*}
        \Big(p_{i_1}\cdots p_{i_r}[M^{4j_1}\times\cdots\times M^{4j_s}]\Big)
    \end{equation*}
    is nonsingular, where both $I=i_1,\dots,i_r$ and
    $J=j_1,\dots,j_s$ range over all partitions of $n$.
\end{theorem}
\begin{proof}
    We only prove the first theorem, 
    the second one is completely analogous.

    Since the basis $c_{i_1}\cdots c_{i_r}$ isn't well behaved
    under the product operation, we change to the basis $s_I(c)$.
    By the product formula, 
    \begin{equation*}
        s_I[K^{j_1}\times\cdots\times K^{j_s}] \neq 0
    \end{equation*}
    only if $I=i_1,\dots,i_r$ is a refinement of $J=j_1,\dots,j_s$.
    Thus by suitable ordering of both the rows and columns, we can assume that
    the matrix is upper triangular. The diagonal elements are
    \begin{equation*}
        s_{i_1,\dots,i_r}[K^{i_1}\times\cdots\times K^{i_r}] =
        s_{i_1}[K^{i_1}]\cdots s_{i_r}[K^{i_r}] \neq 0.
    \end{equation*}
    Thus the matrix is nonsingular.
\end{proof}